%% File: chapters/chapter2.tex
%% Author: Y.U.P. (yashparitkar)
%%
%% Description: Getting started: prerequisites and a first end-to-end capture.

\documentclass[../manual.tex]{subfiles}

\begin{document}

\chapter{Getting Started}
\label{ch:started}

\todo{The highest value chapter in the manual: a reader who follows it from top to bottom gets a waveform on screen without opening any other document. No forward references, so anything needed here is stated here, with the full treatment cross-referenced later.}

\section{Prerequisites}
\label{sec:started_prereq}

The full toolchain is tested using following:

\begin{itemize}
    \item FPGA: Microchip PolarFire SoC Discovery Kit (mpfs-disco-kit)
    \item Libero SoC v2025.2
    \item SoftConsole v2022.2
    \item Ubuntu 22.04.3 LTS
    \item Python 3.10.0
    \item Packages for driver: pyserial (v3.5)
    \item Packages for GUI: PySide6 (v6.5.1)
    \item Rpi Pico with \texttt{debugprobe\_on\_pico} firmware flashed as USB-TTL adapter (other serial adapters will also work)
\end{itemize}

\begin{figure}[H]
    \centering
    \includegraphics[width=0.6\textwidth]{images/replicate_steps/0_testing_setup.jpg}
    \caption{The testing setup.}
    \label{fig:started_setup}
\end{figure}

The tools might or might not work on other versions, but the manual is written for the above. The reader is expected to have a working knowledge of Linux and Python, and to be able to install software and run commands in a terminal.

Not all of it is needed for all of it. Table~\ref{tab:started_prereq} says what each piece buys; leave a row out and everything that does not depend on it still works.

\begin{table}[H]
    \centering
    \begin{tabularx}{\textwidth}{|l|X|}
        \hline
        \rowcolor{gray!40} \textbf{Tool} & \textbf{Needed for} \\
        \hline
        Python 3.10 or later & everything on the host \\
        \hline
        pyserial & the python driver, both front ends \\
        \hline
        PySide6 & the graphical front end, \texttt{--gui}, only \\
        \hline
        USB-TTL adapter & reaching \texttt{libre\_ila\_uart} from the PC \\
        \hline
        Libero SoC & building a core into a design and programming the FPGA \\
        \hline
        SoftConsole & the baremetal C driver, and the demo application this chapter runs \\
        \hline
        GHDL & optional, runs the simulations under \texttt{tests/hdl/} \\
        \hline
        GTKWave & optional, opens the captured \texttt{.vcd}; any VCD viewer will do \\
        \hline
    \end{tabularx}
    \caption{What each prerequisite is for.}
    \label{tab:started_prereq}
\end{table}

\subsection{Installing the host packages}
\label{sec:started_hostpkgs}

Both python packages are in apt on Debian and Ubuntu:

\begin{lstlisting}
sudo apt install python3-serial
sudo apt install python3-pyside6.qtwidgets python3-pyside6.qtgui
\end{lstlisting}

PySide6 ships one package per Qt module there, so \texttt{python3-pyside6} is not a package and \texttt{python3-pyside6.qtcore} on its own is not enough. Everything except \texttt{--gui} works without either of the two.

Ubuntu 22.04 predates PySide6's apt packaging, so on the tested system those two packages do not exist at all. Install it into a virtual environment instead, as in Section~\ref{sec:started_venv}.

The optional tools are one line as well:

\begin{lstlisting}
sudo apt install ghdl gtkwave
\end{lstlisting}

Neither is needed to build a core or take a capture. GHDL runs the simulations under \texttt{tests/hdl/}, and GTKWave is what the driver's \texttt{--waveform-viewer} defaults to when it opens a capture for you.

Finally, reading a serial port on Ubuntu needs membership of the \texttt{dialout} group. Add yourself once and log out and back in, otherwise every connection fails with a permission error on \texttt{/dev/ttyUSB0}:

\begin{lstlisting}
sudo usermod -aG dialout $USER
\end{lstlisting}

\subsection{Creating a Python virtual environment}
\label{sec:started_venv}

This is how PySide6 is installed on a distribution that does not package it, and it is the tidier route on any distribution, since nothing lands system wide.

\begin{enumerate}
    \item Install Python 3.10 or later, pip and the venv module:
    \begin{lstlisting}
sudo apt install python3 python3-pip python3-venv
    \end{lstlisting}
    \item In the root of the repository, create a virtual environment named \texttt{venv-libreila}:
    \begin{lstlisting}
python3 -m venv venv-libreila
    \end{lstlisting}
    \item Activate the virtual environment:
    \begin{lstlisting}
source venv-libreila/bin/activate
    \end{lstlisting}
    \item Install the required packages:
    \begin{lstlisting}
pip install pyserial PySide6
    \end{lstlisting}
\end{enumerate}

The name and the place matter. The root \texttt{Makefile} and the GUI test under \texttt{tests/python/05\_gui/} both look for \texttt{venv-libreila} at the repository root and fall back to it when the system \texttt{python3} has no PySide6, so a virtual environment put anywhere else has to be activated by hand for every command in this manual.

To check that one of the two has it, without running anything else, from the repository root:

\begin{lstlisting}
make check-pyside6
\end{lstlisting}

It prints an install line if neither the system \texttt{python3} nor the virtual environment has PySide6, and stays silent when one of them does.

\section{Getting the sources}
\label{sec:started_sources}

\todo{The clone, then a directory tour in one short table: \texttt{hdl/}, \texttt{codegen/}, \texttt{drivers/}, \texttt{tests/}, \texttt{docs/}.}

\subsection{Cloning the repository}
Make a local clone of the repository, 

\begin{lstlisting}
    git clone https://github.com/yashparitkar/LibreILA
\end{lstlisting}

\subsection{Importing the Libero project}

\begin{enumerate}
    \item Open Libero
    \item In the Project tab $\to$ Open Project $\to$ Browse to \texttt{\{repo\}/tests/full/libre\_ila\_uart/libre\_ila\_uart\_libero/libre\_ila\_uart\_libero.prjx}
    \item You can view the block design in \texttt{Design Heirarchy}, it contains the libre\_ila, axi4smaster\_pktgen (to generate dummy axistream traffic), a simple heartbeat ip and other MPFS blocks.
    \begin{figure}[H]
        \centering
        \includegraphics[width=\textwidth]{images/importing_libero_project/1_exploring_the_bd.png}
        \caption{Exploring the block design.}
        \label{fig:started_libero_bd}
    \end{figure}
    \item Connect the Microchip PolarFire SoC development board to your PC via USB.
    \item In the \texttt{Design Flow} tab, click \texttt{Run PROGRAM Action}, it should do all the synthesis, placement, routing and programming of the FPGA.
    \begin{figure}[H]
        \centering
        \includegraphics[width=\textwidth]{images/importing_libero_project/2_run_program.png}
        \caption{Run PROGRAM Action.}
        \label{fig:started_libero_program}
    \end{figure}
\end{enumerate}

\subsection{Importing the SoftConsole project}
\begin{enumerate}
    \item Open SoftConsole
    \item Right click in the Project Explorer tab $\to$ Import $\to$ General $\to$ Existing Projects into Workspace $\to$ Next $\to$ Browse $\to$ Select \texttt{\{repo\}/tests/full/libre\_ila\_uart/libre\_ila\_uart\_sc} $\to$ Finish
    \begin{figure}[H]
        \centering
        \includegraphics[width=\textwidth]{images/importing_sc_project/1_import_project.png}
        \caption{Importing the project.}
        \label{fig:started_sc_import}
    \end{figure}
    \item You can view the source code in \texttt{src} folder, the file \texttt{src/application/hart1/u54\_1.c} contains the main function which initializes the UART and sends some dummy data over AXI4Stream to the libre\_ila core.
    \begin{figure}[H]
        \centering
        \includegraphics[width=\textwidth]{images/replicate_steps/1_open_the_softconsole_project.png}
        \caption{Opening the SoftConsole project.}
        \label{fig:started_sc_open}
    \end{figure}
    \item Change the Build Configuration to LIM-Debug-DiscoveryKit
    \item Build the project by clicking on the hammer icon in the toolbar.
    \item Debug the project by clicking on the bug icon in the toolbar, it should program the SoC and start the application.
    \begin{figure}[H]
        \centering
        \includegraphics[width=\textwidth]{images/replicate_steps/2_debug_view.png}
        \caption{The Debug view.}
        \label{fig:started_sc_debug}
    \end{figure}
    \item Once the DDR training is complete, the application will stop at breakpoint in \texttt{src/application/hart0/u54\_0.c} file, click on the Resume button in the toolbar to continue the execution of the application. 
\end{enumerate}


\section{Your first capture}
\label{sec:started_first}

\todo{The whole loop on the stock 64-bit AXI4Stream build, every command shown and every output pasted: generate the core, add it to a project, build, connect the USB-TTL adapter, add the device, set a trigger, arm, wait, read back, open the VCD. The trigger is the datasheet's worked example, TVALID and TLAST rising, so the two documents agree.}

\begin{enumerate}
    \item Start the debug project on SoftConsole
    \item Open a the serial terminal to connect with the MSS UART, you can use any serial terminal of your choice, we will use \texttt{minicom} here. Note that, if the drivers for the MPFS is installed properly, you will see three serial devices in \texttt{/dev/} directory, one of them will be streaming the output from the MSS UART. In our case, it is \texttt{/dev/ttyUSB2}.
    \begin{lstlisting}
        minicom -D /dev/ttyUSB2 -b 115200
    \end{lstlisting}
    \item In the serial terminal, you should see the following output:
    \begin{figure}[H]
        \centering
        \includegraphics[width=\textwidth]{images/replicate_steps/4_mss_connected.png}
        \caption{MSS UART connected.}
        \label{fig:started_mss}
    \end{figure}
    \item Configure the pktgen using the menu options, we are using following configuration:
    \begin{lstlisting}
        Packet size: 250 frames
        Packet count: 250 packets
        Packet density: bubble per 2 frames
    \end{lstlisting}
    \begin{figure}[H]
        \centering
        \includegraphics[width=\textwidth]{images/replicate_steps/5_configuring_pktgen.png}
        \caption{Configuring PKTGEN.}
        \label{fig:started_pktgen}
    \end{figure}
    \item Open the GUI front end, you can use the following command to open the GUI:
    \begin{lstlisting}
        python3 libreila.py --gui
    \end{lstlisting}
    \begin{figure}[H]
        \centering
        \includegraphics[width=\textwidth]{images/replicate_steps/3_launch_libre_ila.png}
        \caption{Launching LibreILA.}
        \label{fig:started_launch}
    \end{figure}
    \item In the GUI, click on \texttt{Add Device} and select the serial port to which the USB-TTL adapter is connected, in our case it is \texttt{/dev/ttyACM0}. Click on \texttt{Connect} to connect to the device.
    \item Set the trigger edge, rising and move the position to somewhere middle and trigger condition to \texttt{TVALID and TLAST rising}. Click on \texttt{Arm} to arm the trigger.
    \begin{figure}[H]
        \centering
        \includegraphics[width=\textwidth]{images/replicate_steps/6_configuring_and_arming_the_ila.png}
        \caption{Configuring and arming the ILA.}
        \label{fig:started_arm}
    \end{figure}
    \item Fire the packet generator by selection option \texttt{1. Start PKTGEN transfer} on the MSS serial terminal
    \item The ILA should've captured the data and the ILA status will change to \texttt{Triggered} and then to \texttt{Done} quickly.
    \begin{figure}[H]
        \centering
        \includegraphics[width=\textwidth]{images/replicate_steps/7_firing_pktgen_done.png}
        \caption{Firing PKTGEN, capture done.}
        \label{fig:started_done}
    \end{figure}
    \item Press \texttt{Save} in the \texttt{Capture} section to read data from the ILA and save it to a VCD file.
    \item Press \texttt{Open with gtkwave} to open the capture in gtkwave
    \item In gtkwave, right click on the \texttt{libre\_ila} $\to$ Recure Import $\to$ Insert to add signals to waveform window
    \item Tip: click on the "Magnifying glass with green tick" to fit the capture in the window
    \begin{figure}[H]
        \centering
        \includegraphics[width=\textwidth]{images/replicate_steps/8_gtkwave_window.png}
        \caption{The GTKWave window.}
        \label{fig:started_gtkwave}
    \end{figure}
\end{enumerate}

\section{What just happened}
\label{sec:started_recap}

\todo{Half a page retracing the steps and naming the chapter that covers each in full: the map into the rest of the manual.}
\end{document}
